\documentclass{article}
\usepackage[T1]{fontenc}
\usepackage[utf8]{inputenc}
\usepackage[brazil]{babel}
\usepackage{ismir} % sem "submission": mostra o autor e remove a numeração de linhas
\usepackage{amsmath,amssymb,cite,url}
\usepackage{graphicx}
\usepackage{color}
\usepackage{booktabs}
\usepackage{tikz}

\newcommand{\repourl}{\url{https://github.com/guilhermenovais/boombap-generation}}

\newcommand{\nt}[1]{\textit{#1}}
\newcommand{\tm}[1]{\texttt{#1}}

\title{Geração de Batidas Boom Bap com uma Gramática Probabilística Livre de Contexto}

\oneauthor
  {Guilherme Novais de Souza}
  {Departamento de Ciência da Computação, Universidade Federal de Minas Gerais}

\def\authorname{G. N. de Souza}

\makeatletter
\renewcommand{\permission}{%
  \begin{figure}[b]{%
    \footnotesize
    \begin{spacing}{1.1}
      \copyright\ \authorname. Trabalho prático~1 de DCC831 -- Inteligência Artificial na Música,
      DCC/UFMG, 2026. Relatório no formato ISMIR.
    \end{spacing}
    }
  \end{figure}
}
% Rótulos que o ismir.sty fixa em inglês
\def\abstract{\begin{center}\bf RESUMO\end{center}}
\def\fnum@figure{\textbf{Figura \thefigure}}
\def\fnum@table{\textbf{Tabela \thetable}}
\def\thebibliography#1{%
  \section{Referências}\list
  {[\arabic{enumi}]}{%
  \settowidth\labelwidth{[#1]}\leftmargin 1em
  \advance\leftmargin\labelsep
  \usecounter{enumi}%
  }%
  \def\newblock{\hskip .01em plus .01em minus .01em}%
  \sloppy\clubpenalty4000\widowpenalty4000
  \sfcode`\.=1000\relax
}
\makeatother

\sloppy

\begin{document}

\maketitle

\begin{abstract}
Este trabalho apresenta um gerador de batidas de boom bap, o hip-hop dos anos 1990, baseado só na
expansão de uma gramática probabilística livre de contexto (PCFG): uma parte estrutural indexada deriva
a forma da música com o número exato de compassos, e uma tabela de regras define o estilo. A
amostragem é reprodutível pela semente e monotônica nos hiperparâmetros, e em 100 sementes sorteadas
nenhuma música violou as restrições de estilo.
\end{abstract}

\section{Introdução}\label{sec:intro}

Gramáticas formais são uma representação natural para músicas com forma hierárquica e
convencionada~\cite{Roads79,Lerdahl83}. Uma batida de hip-hop é uma árvore: a música se divide em
seções, as seções em blocos de 4 ou 8 compassos, e os compassos em eventos de percussão numa grade.
O boom bap tem groove padronizado (bumbo no primeiro tempo, caixa nos
tempos 2 e 4, swing de MPC), mas varia dentro do padrão. Este trabalho descreve o \tm{boombap}, que gera músicas completas nesse estilo, em MIDI, só
pela expansão de uma PCFG em que (i)~forma e duração são garantidas por construção, (ii)~as restrições
de estilo são verificadas sobre a tabela de regras antes da geração e (iii)~a amostragem é
reprodutível e monotônica nos hiperparâmetros.

\section{Trabalhos Relacionados}\label{sec:related}

Roads e Wieneke~\cite{Roads79} fazem um panorama do uso de gramáticas como representação musical. O
trabalho mais próximo deste é o Bol Processor de Kippen e Bel~\cite{Kippen92}, em que gramáticas
formais representam sequências de toques de tabla: como aqui, os terminais são eventos percussivos e a
gramática codifica um estilo específico. Keller e Morrison~\cite{Keller07} usam uma gramática
probabilística para improvisar jazz, com terminais abstratos que só depois viram notas, separação
parecida com a nossa entre derivação e realização.

\section{Método}\label{sec:method}

\subsection{Restrição de Estilo e Representação}

O estilo alvo é o boom bap dos anos 1990: compasso 4/4, 85--95~bpm, swing de MPC, forma
intro~$\to$ verso~$\to$ refrão~$\to$ verso~$\to$ outro, com cada seção composta de blocos de 4 ou 8
compassos. Cada compasso é uma grade de 16 semicolcheias (posições 0--15, com 120 ticks cada e
480 ticks por semínima). Um terminal de bateria tem a forma \tm{peça@posição:velocity}, por exemplo
\tm{kick@0:118}, mapeado para o canal 10 do General MIDI. Os outros terminais são marcadores: \tm{grid8}/\tm{grid16}
dá a grade do chimbal (e quais posições recebem swing), \tm{dv$\pm$k} e \tm{dt$\pm$k} são desvios de
velocity e de tempo, \tm{$\langle$seção:X$\rangle$} marca uma seção, e \tm{deg$k$} é um grau da escala
menor natural, tocado pelo baixo a cada ataque de bumbo, com a tônica na oitava 2.

\subsection{A Gramática}

A gramática é $G=(N,\Sigma,R,S)$ com uma função $P:R\to[0,1]$ tal que
$\sum_{\alpha} P(A\to\alpha)=1$ para todo $A\in N$. O axioma é $S=\nt{Música}(n)$, com $n$ o
número de compassos pedido. Para $n=40$, a listagem do programa tem 193 não-terminais,
100 terminais e 386 regras. Destes, 121 têm duas ou mais produções com $0<P<1$, de modo que a gramática é
não-determinística: o mesmo axioma gera músicas diferentes.

\begin{figure*}[t]
\footnotesize
\centering
\begin{tabular}{@{}l@{\;}c@{\;}l@{\quad}l@{}}
\toprule
\nt{Música}$(n)$ & $\to$ & \nt{Forma}$(n)$ & $[1]$ \\
\nt{Forma}$(n)$ & $\to$ & \nt{Intro}$(i)$ \nt{Corpo}$(n{-}i{-}o)$ \nt{Outro}$(o)$ & $[\propto P_I(i)\,P_O(o)]$ \\
\nt{Corpo}$(m)$ & $\to$ & \nt{Verso}$(a)$ \nt{Refrão}$(b)$ \nt{Verso}$(c)$ & $[\propto w(a,b,c)]$ \\
\nt{Verso}$(k)$ & $\to$ & \tm{$\langle$seção:verso$\rangle$} \nt{Seção}$(k,g,\textrm{início\_seção})$ & $[P_G(g)]$ \\
\nt{Seção}$(k,g,s)$ & $\to$ & \nt{Bloco4}$(g,s)$ \nt{FimBloco}$(k{-}4,g)$ & $[0{,}4]$ \\
 & $\mid$ & \nt{Bloco8}$(g,s)$ \nt{FimBloco}$(k{-}8,g)$ & $[0{,}6]$ \\
\nt{Bloco4}$(g,s)$ & $\to$ & \nt{Compasso}$(g,\textrm{primeiro},s)$ \nt{Compasso}$(g,\textrm{meio},\textrm{normal})^{2}$ & $[1]$ \\
\nt{FimBloco}$(k,g)$ & $\to$ & \nt{Compasso}$(g,\textrm{fim},\textrm{virada})$ \nt{Seção}$(k,g,\textrm{pós\_virada})$ & $[\textrm{@fill\_end}=0{,}30+0{,}70f]$ \\
 & $\mid$ & \nt{Compasso}$(g,\textrm{fim},\textrm{groove})$ \nt{Seção}$(k,g,\textrm{normal})$ & [resto] \\
\nt{Compasso}$(g,\cdot,s)$ & $\to$ & \nt{Prato}$(s)$ \nt{Fill\_}$g$ $\mid$ \nt{Prato}$(s)$ \nt{Groove\_}$g$ & $[\textrm{@fill\_mid}=0{,}05+0{,}25f \mid \textrm{resto}]$ \\
\nt{Prato}$(s)$ & $\to$ & \tm{crash@0:110} $\mid$ $\varepsilon$ & $[p_s \mid \textrm{resto}]$ \\
\midrule
\nt{Groove\_Classic} & $\to$ & \tm{kick@0:118 snare@4:115 kick@10:104 snare@12:117} & \\
 & & \nt{OptKicks\_Classic} \nt{Ghosts\_Classic} \nt{Hats\_Classic} & $[1]$ \\
\nt{OptKick7\_Classic} & $\to$ & \tm{kick@7:88} $\mid$ $\varepsilon$ & $[0{,}20+0{,}55d \mid \textrm{resto}]$ \\
\nt{Hats\_Classic} & $\to$ & \nt{Hats16} $\mid$ \nt{Hats8\_Classic} & $[0{,}10+0{,}50d \mid \textrm{resto}]$ \\
\bottomrule
\end{tabular}
\caption{Acima, as famílias estruturais indexadas; abaixo, algumas regras da tabela de estilo
\tm{boombap.rules}. $d$ e $f$ são \tm{density} e \tm{fills}. Em \nt{Forma}, $i,o\in\{4,8\}$; em
\nt{Corpo}, $a,b,c$ são múltiplos de 4 com $a{+}b{+}c=m$ e
$w=\exp(-\lVert (a,b,c)/m-(0{,}3;\,0{,}4;\,0{,}3)\rVert^2/2\sigma^2)$, $\sigma=0{,}08$;
$g\in\{$Classic, Lazy, Bounce$\}$ com $P_G=(0{,}45;\,0{,}25;\,0{,}30)$. \nt{Intro}, \nt{Refrão} e \nt{Outro} são
análogas a \nt{Verso}; $\nt{Seção}(0,g,s)\to\varepsilon$.}
\label{fig:grammar}
\end{figure*}

\textbf{Parte estrutural.} As famílias da Figura~\ref{fig:grammar} (acima) são geradas pelo código,
mas todos os seus pesos vêm de parâmetros da tabela de regras. Os índices são o número de compassos
$k$, o groove $g$ e um estado $s$. Toda regra conserva a soma dos índices de duração do lado direito,
logo qualquer derivação de $\nt{Música}(n)$ tem exatamente $n$ compassos, sem amostragem por
rejeição. O groove é sorteado uma vez por seção, e o estado $s$ registra o que precede o compasso (por
exemplo, um prato após uma virada). Para cada $n$ os índices tomam valores finitos; instanciá-los dá
um conjunto finito de não-terminais comuns, e a gramática continua livre de contexto.

\textbf{Parte de estilo.} A tabela \tm{boombap.rules} define três grooves (Classic, Lazy e Bounce),
com bumbos e caixas fixos, notas opcionais e um padrão de chimbal; viradas \nt{Fill\_}$g$, que trocam
o final do compasso por uma de cinco ``caudas'' (rufo, bumbo--caixa etc.); e a humanização: cada
terminal de bateria $x$ equivale a $X\to x\;\nt{HumV}\;\nt{HumT}$. Com o baixo ligado, cada bumbo
também expande \nt{BassFirst} ou \nt{BassDeg}, que sorteiam o grau.

\textbf{Restrições validadas.} Antes de cada geração, a tabela é verificada em toda a grade de
\tm{density} e \tm{fills}: probabilidades somam 1; todo groove tem bumbo no 0, caixa no 4 e no 12 e um
bumbo sincopado; toda virada mantém bumbo no 0 e caixa no 4; todo compasso tem chimbal e uma grade;
\tm{@fill\_end}~$>$~\tm{@fill\_mid}; e não há recursão infinita. Senão, a tabela é recusada.

\subsection{Realização: dos Terminais ao MIDI}

Um evento na posição $p$ do compasso $c$ toca no tick $1920c+120p+\delta(p)+\textit{dt}$, limitado ao
compasso. No swing da MPC,
com $w$\%, as posições fracas da grade atrasam $\delta=(w/100-1/2)\cdot 480$ ticks nas colcheias (metade
nas semicolcheias): $w=50$ é reto e $w\approx 67$ é ternário. \tm{dv} soma-se à velocity. O MIDI
guarda a semente e os hiperparâmetros, e a derivação completa vai para um \tm{.json}.

\textbf{Hiperparâmetros.} \tm{density} e \tm{fills} (0--1, padrão 0,5) controlam notas opcionais e
viradas e entram como expressões nas probabilidades da tabela, como $0{,}20+0{,}55\,d$ em
\nt{OptKick7\_Classic}. Os demais são \tm{swing} (52--75\%), \tm{bpm} (85--95), \tm{bars} (20--256, múltiplo de 4),
\tm{bass}/\tm{key} (baixo em tom menor) e a semente.

\subsection{Reprodutibilidade}

O número $u\in[0,1)$ de cada nó vem do hash
$\textrm{sha256}(\textit{semente}\,|\,\textit{caminho})$ do seu caminho na árvore, e a regra é escolhida pela CDF inversa na ordem da tabela. A mesma semente gera arquivos
idênticos, e mudar um parâmetro não muda nenhum $u$. Como as probabilidades são não-decrescentes no parâmetro e a alternativa ``presente'' vem
primeiro, $u<p(d_1)\Rightarrow u<p(d_2)$ para $d_1\le d_2$: toda nota opcional presente com
\tm{density} baixa continua presente com \tm{density} alta.

\section{Resultados}\label{sec:results}

Geramos quatro amostras (Tabela~\ref{tab:results}), cada uma com \tm{.mid}, \tm{.json} e \tm{.wav}
(FluidSynth): \tm{seed7\_default} (padrões),
\tm{seed7\_busy} (mesma semente, $d=f=1$), \tm{seed21\_bass\_dm} (outra semente, swing 70 e baixo em
Ré menor) e \tm{seed99\_long} (86~bpm, 64 compassos, $d=0{,}2$, $f=0{,}3$). Os arquivos estão em
\tm{samples/} no repositório do projeto\footnote{\repourl}.

\begin{table}[t]
\footnotesize
\centering
\setlength{\tabcolsep}{3.4pt}
\begin{tabular}{@{}lrrrrrr@{}}
\toprule
Amostra & bumbo & caixa & aro & chimbal & viradas & atraso \\
\midrule
\tm{seed7\_default} & 2,52 & 2,48 & 0,55 & 8,55 & 10/40 & 33,1 \\
\tm{seed7\_busy} & 2,77 & 2,75 & 0,75 & 8,85 & 17/40 & 31,1 \\
\tm{seed21\_bass\_dm} & 2,42 & 2,88 & 0,28 & 9,55 & 10/40 & 74,6 \\
\tm{seed99\_long} & 2,50 & 2,56 & 0,14 & 8,86 & 13/64 & 33,6 \\
\midrule
seed 7, $d=0$ & 2,35 & 2,33 & 0,28 & 7,30 & 10/40 & 37,8 \\
seed 7, $d=1$ & 2,83 & 2,52 & 0,85 & 9,43 & 10/40 & 30,3 \\
\bottomrule
\end{tabular}
\caption{Saída de \tm{boombap analyze}: notas por compasso (chimbal fechado), compassos com virada e
atraso médio das posições fracas, em ticks (120 = uma semicolcheia).}
\label{tab:results}
\end{table}

Entre \tm{seed7\_default} e \tm{seed7\_busy} a forma é a mesma (4--12--12--8--4), porque nenhuma regra
estrutural depende de \tm{density}. As viradas passam de 10 para 17 compassos e os bumbos, caixas e
aros aumentam. Variando só \tm{density} (duas últimas linhas), o número de viradas fica igual e
bumbo, aro e chimbal crescem. Com swing 70,
\tm{seed21\_bass\_dm} tem mais que o dobro do atraso das demais.



\textbf{Validação.} Os testes geram 100 sementes com hiperparâmetros sorteados e analisam cada MIDI:
nenhuma violou as restrições de estilo, e as seções somam o número de compassos pedido. Regenerar com os
parâmetros gravados reproduz o MIDI byte a byte. A monotonicidade em \tm{density} vale nas 10 sementes testadas, mas em 3 delas \tm{density} 0$\to$1 multiplicou bumbo+chimbal por menos que a meta
de 1,3 (1,27 a 1,29): os coeficientes de \tm{density} poderiam ser maiores.

\section{Discussão e Trabalhos Futuros}\label{sec:discussion}

O sistema gera só bateria e baixo, sem melodia, harmonia ou samples. As probabilidades foram escritas à
mão, e, por ser livre de contexto, a gramática não repete motivos entre seções; também não houve teste
com ouvintes. Como trabalhos futuros, as probabilidades poderiam ser estimadas do Groove MIDI Dataset~\cite{Gillick19} e
gramáticas sensíveis ao contexto permitiriam repetir motivos.

\bibliography{refs}

\end{document}
